\chapter{Coexistence and Gradual Migration Strategy}

The transition from the legacy SRP system to the new ServiceRequestManager architecture represents a substantial technical and operational challenge that requires careful planning to minimize business disruption while ensuring data integrity and system reliability. Rather than implementing a disruptive "big bang" migration that would introduce significant risk and potential service interruptions, the project adopts a phased coexistence strategy that enables gradual migration of customers and service request types while maintaining full operational continuity.

\section{Migration Strategy Overview}

The migration architecture implements a dual-flow execution model where both SRP and SRM systems operate simultaneously, with individual service request actions configured to route through either the legacy or modern execution pipeline based on the \texttt{target\_type} attribute of the Action model. This attribute, defined as an enumeration in \texttt{ActionTargetTypes}, determines the execution mechanism for each action:

\begin{minted}[breaklines,fontsize=\small]{python}
class ActionTargetTypes(Enum):
    DOMAIN_DC = "From SR Domain (Domain Controller)"
    DOMAIN_EXC = "From SR Domain (Exchange Server)"
    INTERNAL_WORKER = "Internal Worker"
    INTERNAL_API = "API Call"
    SRP = "SRP Agent (Domain Controller)"
\end{minted}

Actions configured with \texttt{target\_type = SRP} continue routing through the legacy SRP Agent execution mechanism, preserving the existing workflow without modification. Actions configured with any other target type utilize the new Ansible-based execution pipeline through Provision Prime, benefiting from improved logging, error handling, and state management capabilities introduced in SRM.

This per-action migration control provides several strategic advantages: it enables pilot deployments on specific customers or action types, allowing validation of the new architecture in production with controlled scope; it maintains business continuity by ensuring that any discovered issues affect only the subset of actions migrated to the new flow; it supports progressive learning, enabling the operations team to develop expertise with the new system incrementally; and it provides fallback capability, as actions can be reverted to SRP execution if critical issues are discovered during migration.

The primary architectural trade-off of this approach is that the catalog configuration must be maintained in the SRP database during the coexistence period, as SRP lacks the sophisticated configuration features introduced in SRM (three-level hierarchy, external API configurations, template-based autofill). The SRM system operates as a read-only consumer of catalog data from SRP, applying its advanced configuration features as an overlay on the imported base catalog.

\section{Catalog Synchronization from SRP to SRM}

The catalog synchronization mechanism implements a unidirectional data flow where SRP serves as the authoritative source of catalog definitions, with SRM periodically importing and enriching this data with its extended configuration capabilities. This synchronization is executed through the \texttt{sync\_srp\_catalog} Django management command, scheduled as a Kubernetes CronJob for daily execution at 01:00 AM.

\subsection{Synchronization Architecture}

The synchronization process establishes a direct ODBC connection to the SRP SQL Server database and executes a series of import operations that reconstruct the complete catalog hierarchy in the SRM PostgreSQL database:

\begin{minted}[breaklines,fontsize=\small]{python}
SRP_CONN_STR = (
    "DRIVER={ODBC Driver 18 for SQL Server};"
    f"SERVER={os.getenv('SRP_SERVER')}\\SRP;"
    f"DATABASE={os.getenv('SRP_DB')};"
    f"UID={os.getenv('SRP_UID')};"
    f"PWD={os.getenv('VAULT_SRP_PASSWORD')};"
    "TrustServerCertificate=yes;"
)

with disable_auditlog(), pyodbc.connect(SRP_CONN_STR) as conn:
    cursor = conn.cursor()
    sync_action_types(cursor)
    sync_action_categories(cursor)
    sync_action_groups(cursor)
    sync_customers(cursor)
    sync_actions(cursor)
    sync_config_parameters(cursor)
    sync_fields(cursor)
    sync_customer_actions(cursor)
\end{minted}

The synchronization executes with audit logging disabled through the \texttt{disable\_auditlog()} context manager to prevent creation of thousands of audit log entries for bulk import operations, improving performance and reducing database bloat. Each synchronization function implements a consistent pattern: 
\begin{enumerate}
    \item Query the source table in SRP;
    \item Construct Django model instances for new or modified records;
    \item Execute bulk create/update operations.
\end{enumerate}

\subsection{Differential Synchronization Logic}

The synchronization functions implement sophisticated differential logic that minimizes database operations by detecting actual changes rather than blindly recreating all data. The \texttt{sync\_action\_categories} function exemplifies this pattern:

\begin{minted}[breaklines,fontsize=\small]{python}
cursor.execute("SELECT * FROM TBCatalog_Categories")

records: List[Dict[str, Any]] = [
    dict(zip(columns, row)) for row in cursor.fetchall()
]

source_ids: Set[int] = {int(r["IDCategory"]) for r in records}

existing_categories: Dict[int, ActionCategory] = {
    ac.id: ac
    for ac in ActionCategory.objects.filter(id__in=source_ids)
}

to_create: List[ActionCategory] = []
to_update: List[ActionCategory] = []

for record in records:
    category_id: int = int(record["IDCategory"])

    if category_id in existing_categories:
        obj: ActionCategory = existing_categories[category_id]
        obj.description = record["DescCategory"]
        obj.display_order = record["OrderN"]
        to_update.append(obj)
    else:
        to_create.append(
            ActionCategory(
                id=category_id,
                description=record["DescCategory"],
                display_order=record["OrderN"],
            )
        )

if to_create:
    ActionCategory.objects.bulk_create(to_create)

if to_update:
    ActionCategory.objects.bulk_update(
        to_update,
        ["description", "display_order"],
    )
\end{minted}

This approach constructs in-memory dictionaries of existing records keyed by primary key, then iterates through source records to segregate them into create and update lists. Bulk operations execute with batch sizes appropriate for PostgreSQL, significantly outperforming individual save() calls. The \texttt{sync\_deletions} helper function identifies records present in SRM but absent from SRP, removing obsolete catalog entries to maintain synchronization.

Specifically, an initial full synchronization strategy was adopted, which involved deleting and recreating all records at each job execution. However, this approach proved to be inefficient and problematic with SRP's large data volumes, resulting in average execution times of around 2 hours. The current differential logic significantly reduces the load on the database, bringing average execution times down to approximately 20 seconds.

\subsection{Relationship Resolution and Foreign Key Management}

Catalog synchronization must handle complex foreign key relationships where referenced entities may not yet exist when dependent records are processed. The synchronization order (types → categories → groups → customers → actions) ensures that parent entities exist before children reference them. For many-to-many relationships and optional foreign keys, the sync functions implement validation that skips records referencing non-existent entities, logging warnings for manual investigation.

The \texttt{sync\_actions} function demonstrates this validation pattern:

\begin{minted}[breaklines,fontsize=\small]{python}
customer_ids: Set[int] = {
    int(r.get("IDCustomer", -9999))
    for r in records
    if r.get("IDCustomer", -9999) != -9999
}

existing_customer_ids: Set[int] = set(
    Customer.objects
        .filter(id__in=customer_ids)
        .values_list("id", flat=True)
)

for record in records:
    customer_id: Optional[int] = record.get("IDCustomer", -9999)

    if customer_id == -9999 or customer_id not in existing_customer_ids:
        customer_id = None  # Action not customer-specific
\end{minted}

This pattern prevents foreign key constraint violations while maintaining referential integrity, with nullable foreign keys set to None when references cannot be resolved.

\subsection{Post-Import Enrichment}

Following the base catalog import from SRP, the synchronization process executes several enrichment operations that apply SRM-specific configurations:

\textbf{API Configuration Synthesis}: The \texttt{sync\_api\_configurations()} function creates \texttt{ExternalAPIConfiguration} records that define option sources for select fields, mapping SRP's simpler parameter-based options to SRM's sophisticated external API integration framework.

\textbf{Field Duplication Resolution}: The \texttt{fix\_duplicated\_field\_names()} function identifies and resolves scenarios where multiple fields with identical names exist within the same action scope (considering action-level and customer-action-level configurations), a situation that would cause ambiguity in field references during form rendering and script parameter generation.

\textbf{HDA Field Configuration}: The \texttt{sync\_hda\_fields()} function synthesizes configuration for HDA external fields and classification fields based on ticket type mappings, creating the metadata required for the frontend to render these additional input controls.

\textbf{Customer-Specific Customizations}: The \texttt{sync\_customer\_specific\_customizations()} function processes SRP-specific configuration tables and translates them into SRM's three-level configuration hierarchy, creating \texttt{CustomerAction} and \texttt{CustomerActionField} records that override default behavior.

\subsection{Sequence Reset and Consistency Verification}

The synchronization concludes by resetting PostgreSQL auto-increment sequences to prevent primary key conflicts on subsequent manual record creation:

\begin{minted}[breaklines,fontsize=\small]{python}
def reset_sequences(app_label: str):
    for model in apps.get_app_config(app_label).get_models():
        cursor.execute(f"""
            SELECT setval(
                pg_get_serial_sequence(
                    '{model._meta.db_table}', 'id'
                ),
                COALESCE(MAX(id), 1)
            ) FROM {model._meta.db_table};
        """)
\end{minted}

This query identifies the maximum existing primary key for each table and adjusts the sequence generator accordingly, ensuring that future inserts receive non-conflicting IDs.

\section{Execution History Synchronization}

While catalog synchronization flows unidirectionally from SRP to SRM, execution history must be merged from both systems to provide unified visibility into all service requests regardless of which execution pipeline processed them. The \texttt{sync\_srp\_sr\_history} command implements this bidirectional synchronization, scheduled to execute at 02:00 AM (one hour after catalog sync to ensure current configurations are available).

\subsection{Service Request Import Architecture}

The execution history synchronization imports three primary entity types: service request headers, field values, and execution logs. The synchronization is filtered by a configurable date threshold to avoid importing the complete historical archive:

\begin{minted}[breaklines,fontsize=\small]{python}
start_date: str = os.getenv(
    "SRP_SR_SYNC_START_DATE",
    datetime.date.today().replace(day=1).strftime("%Y-%m-%d")
)
cursor.execute(
    f"SELECT * FROM TBSR_Headers WHERE Date > '{start_date}'"
)
\end{minted}

This date-based filtering defaults to the first day of the current month, importing only recent execution history while avoiding performance degradation from processing years of historical data. The threshold is configurable through environment variables to support one-time historical imports or adjusted synchronization windows.

\subsection{Field Value Encryption and Sensitive Data Handling}

Service request field values imported from SRP must respect SRM's enhanced security model where sensitive fields (passwords, credentials) are encrypted at rest. The \texttt{sync\_service\_request\_fields} function implements transparent encryption during import:

\begin{minted}[breaklines,fontsize=\small]{python}
for record in records:
    field_id: int = int(record["IDField"])
    value: str = record["Value"]
    field: Field = fields_map.get(field_id)
    
    if field and field.is_sensitive:
        value: str = encrypt_sensitive_data(plaintext=value)
    
    ServiceRequestField.objects.create(
        service_request_id=service_request_id,
        field_id=field_id,
        value=value
    )
\end{minted}

This encryption occurs transparently during synchronization, ensuring that legacy SRP data benefits from SRM's enhanced security protections without requiring modifications to the source system.

\subsection{Log Unification and Type Mapping}

Execution logs stored in SRP use a numeric type code system incompatible with SRM's enumerated log types. The synchronization implements a mapping table that translates SRP type codes to SRM equivalents:

\begin{minted}[breaklines,fontsize=\small]{python}
log_type_mapping = {
    "12": "0",  # Creation
    "5": "1",   # Approval
    "9": "2",   # Execution Started
    "8": "3",   # Execution Completed
    # ... additional mappings
}

for record in records:
    log_type: str = record["TipoLog"]
    if log_type:
        log_type = log_type_mapping.get(str(log_type))
\end{minted}

This mapping preserves log semantics across systems, enabling unified filtering and analysis of execution events regardless of their originating system.

\section{Dual Execution Flow Management}

The coexistence architecture's core mechanism is the dynamic routing of service request execution based on action configuration. When a service request is created through the SRM API, the system evaluates the action's \texttt{target\_type} to determine whether to invoke the new or legacy execution pipeline.

\subsection{Legacy Flow Invocation}

For actions configured with \texttt{target\_type = SRP}, the service request creation serializer implements specialized handling that creates the HDA ticket but delegates service request entity creation to SRP:

\begin{minted}[breaklines,fontsize=\small]{python}
is_srp_flow: bool = action.target_type == ActionTargetTypes.SRP.name

# Create HDA ticket
response_ticket: dict = mysupport.create_ticket(
    codicli=codicli, idrs=customer.idrs, body=body
)
ticket_id: str = response_ticket.get("ID")

if is_srp_flow:
    logger.info(
        f"SRP flow detected. Ticket {ticket_id} created. "
        f"ServiceRequest will be imported by cronjob."
    )
    return ServiceRequest(
        action=action, 
        customer=customer, 
        ticket_reference=ticket_id
    )
\end{minted}

This pattern creates the ticket through SRM's HDA integration (ensuring consistent ticket creation regardless of execution flow), but returns a non-persisted ServiceRequest object to the API consumer. The actual service request record will be created by SRP's legacy REST API when HDA invokes it with the ticket reference, then synchronized into SRM's database by the nightly \texttt{sync\_srp\_sr\_history} job.

\subsection{Modern Flow Execution}

For actions configured with any non-SRP target type, the service request creation proceeds through the complete SRM pipeline: the ServiceRequest record is persisted immediately with NEW status, the ticket reference is stored, and the asynchronous Celery task chain is initiated to orchestrate execution through Provision Prime and Ansible. This execution follows the complete workflow described in previous chapters, including validation, state transitions, approval routing, and execution monitoring.

\subsection{Transparent Frontend Integration}

The dual-flow architecture remains transparent to frontend applications, which interact exclusively with SRM APIs regardless of the underlying execution mechanism. The service request creation endpoint (\texttt{POST /api/customers/\{IDRS\}/service-requests/}) accepts identical payloads for both SRP and SRM actions, with the backend handling execution routing internally. This abstraction ensures that frontend applications require no modifications as actions are gradually migrated from SRP to SRM execution.

\section{Pilot Strategy and Progressive Rollout}

The migration strategy implements a risk-managed progression that validates the new architecture incrementally before full deployment, minimizing exposure to potential issues while building operational confidence.

\subsection{Phase 1: Internal Pilot}

The initial migration phase configures a small subset of actions used exclusively by internal IT operations teams to execute through the SRM pipeline. These actions are selected based on criteria including low execution volume (limiting impact scope), non-critical business processes (tolerating potential issues), and comprehensive logging requirements (benefiting from SRM's enhanced observability). Internal teams provide direct feedback on execution behavior, error handling, and operational experience, enabling rapid iteration on issues discovered during real-world usage.

\subsection{Phase 2: Customer Pilot}

Following successful internal validation, the migration expands to a carefully selected cohort of pilot customers. Pilot customers are chosen based on criteria including technical sophistication (ability to articulate detailed issue reports), diverse action usage patterns (validating various workflow types), and strong relationship with the company (tolerance for minor issues during pilot phase). Each pilot customer has a subset of their actions migrated to SRM execution, typically starting with low-risk operations like information queries before progressing to more critical provisioning workflows.

The pilot phase implements enhanced monitoring including dedicated Grafana dashboards tracking pilot action execution metrics, daily review meetings to assess execution success rates and identify issues, and direct communication channels with pilot customers for immediate feedback. Issues discovered during pilot are addressed through hotfix deployments, with actions reverted to SRP execution if critical problems cannot be resolved quickly.

\subsection{Phase 3: Progressive Expansion}

Successful pilot completion triggers progressive expansion where additional customer cohorts are migrated on a weekly cadence. Each expansion wave includes a defined set of customers and actions, with execution metrics monitored for anomalies before proceeding to the next wave. The expansion velocity is dynamically adjusted based on observed success rates: high success rates (>99\% execution completion without errors) enable acceleration, while elevated error rates trigger pause for investigation and remediation.

The rollout maintains a progressive migration of action types from simple to complex: read-only information queries migrate first, followed by non-destructive provisioning operations (user creation, group membership addition), then critical operations (password resets, account deletions), and finally complex multi-step workflows. This ordering ensures that the most fault-tolerant operations migrate first, building confidence before tackling higher-risk scenarios.

\subsection{Phase 4: SRP Decommissioning}

Migration completion occurs when all customers and actions have successfully transitioned to SRM execution, validated through sustained periods (30+ days) of stable operation without SRP fallback. At this point, the SRP system transitions to read-only mode, maintaining historical data for reference while no longer accepting new service request creation. The catalog synchronization job is deactivated, with SRM becoming the authoritative catalog source. Finally, the SRP infrastructure is decommissioned after a retention period ensuring all audit and compliance requirements are satisfied.

\section{Data Integrity and Consistency Mechanisms}

The coexistence architecture implements several mechanisms to maintain data integrity and consistency across the distributed system during the migration period.

\subsection{Transactional Boundaries}

Synchronization operations execute within database transaction contexts that ensure atomicity: either all records in a synchronization batch commit successfully, or none do, preventing partial updates that would leave the catalog in an inconsistent state. The use of \texttt{bulk\_create} and \texttt{bulk\_update} operations with explicit field lists ensures that only intended modifications occur, avoiding inadvertent overwrites of SRM-specific enhancements applied on top of imported SRP data.

\subsection{Idempotent Synchronization}

The synchronization logic is designed for idempotency, enabling safe re-execution in case of failures. Each sync function can be invoked multiple times with identical results, as record identification is based on deterministic primary keys from SRP. Failed synchronization jobs can be retried without concern for duplicate record creation or data corruption, with the differential logic ensuring that only actual changes trigger database operations.

\subsection{Referential Integrity Validation}

Foreign key relationships are validated before record creation to prevent orphaned references. The synchronization order ensures parent entities exist before children, while optional relationships are nulled when references cannot be resolved. Post-synchronization verification queries identify any referential integrity violations that escaped validation, with alerts triggered for manual investigation and resolution.

\subsection{Audit Trail Preservation}

While audit logging is disabled during bulk synchronization for performance, the synchronization jobs themselves are comprehensively logged through the monitoring infrastructure. Execution start/end times, record counts, errors encountered, and performance metrics are recorded in New Relic, providing complete visibility into synchronization behavior and enabling trend analysis to detect degradation or anomalies.

\section{Challenges and Operational Considerations}

The coexistence strategy introduces several operational challenges and architectural trade-offs that require careful management during the migration period.

\subsection{Dual Maintenance Burden}

The primary operational challenge is the requirement to maintain catalog configurations in the legacy SRP system during coexistence. Any configuration changes (new actions, field modifications, customer onboarding) must be performed in SRP and validated through the synchronization pipeline before becoming effective in SRM. This dual-system maintenance increases operational complexity and the potential for configuration drift if synchronization jobs fail or are delayed.

\subsection{Synchronization Latency}

The nightly synchronization schedule introduces up to 24-hour latency between configuration changes in SRP and their visibility in SRM. For urgent configuration updates required to resolve production issues, manual synchronization execution or direct database modification in SRM may be necessary, introducing risk and complexity. The project mitigates this through careful change management processes and pre-planned configuration updates executed during maintenance windows.

\subsection{Execution History Fragmentation}

During the coexistence period, service request execution history is split between SRP and SRM systems, with the unified view available only in SRM after synchronization. Operational staff must be trained to query SRM exclusively for execution history rather than accessing SRP directly, ensuring consistent data interpretation. The synchronization job's recency filter means that very old historical data may not be available in SRM, requiring retained access to SRP for historical research.

\subsection{Performance Considerations}

The synchronization jobs execute significant database operations (tens of thousands of record comparisons and updates) that consume computational resources and database connections. The jobs are scheduled during low-traffic periods (early morning) to minimize impact on production operations, with timeout configurations adjusted to accommodate long-running synchronization of large datasets. Monitoring ensures that synchronization completion times remain within acceptable bounds (< 30 minutes), with alerts triggered if execution duration increases beyond thresholds indicating performance degradation.

The coexistence and gradual migration strategy provides a pragmatic path to modernizing the service request management infrastructure while maintaining operational continuity and minimizing risk. The dual-flow architecture enables incremental validation of the new system in production environments, building confidence through measured progression rather than high-risk wholesale replacement. While the strategy introduces operational complexity through the need for synchronization and dual-system maintenance, these costs are justified by the substantial risk reduction and business continuity assurance provided by the gradual approach.
